\documentclass{ximera}

\begin{document}

    \title{In-Class Assignment 5}
    
    \begin{abstract} Solve the following problems.
    \end{abstract}
    \maketitle


\begin{exercise}
Consider the matrix
\[
A = \begin{pmatrix}
10^{-4} & 1\\
1 & 1
\end{pmatrix}.
\]

\begin{enumerate}
    \item Perform Gaussian elimination without pivoting to obtain an upper triangular matrix $U$. What is the multiplier used in the elimination step?

    \item Now perform Gaussian elimination with partial pivoting. What row exchange is performed, and what is the resulting upper triangular matrix?

    \item Explain why partial pivoting is preferable in this example. In particular, compare the sizes of the multipliers obtained in the two approaches.
\end{enumerate}
\end{exercise}



\end{document}